% Chapter 6: Discussion
% Target: 6-8 pages (1,500-2,000 words)

\section{Hypothesis Validation}

\subsection{H1: LLM vs ML Performance}

\textbf{Hypothesis}: LLM-based trading strategies achieve superior risk-adjusted returns compared to traditional ML approaches.

\textbf{Finding}: The hypothesis is \textbf{supported}. DeepSeek V3.2 achieved a Sharpe ratio of 0.66 on CSI300, significantly outperforming the best ML model (LSTM with 0.17). On QQQ, DeepSeek V3.2 achieved a Sharpe ratio of 1.83, compared to LightGBM's 1.07.

\textbf{Statistical Evidence}: Jobson-Korkie tests confirm the difference in Sharpe ratios is statistically significant ($p < 0.05$).

\subsection{H2: Risk Control During Market Stress}

\textbf{Hypothesis}: LLM strategies demonstrate better risk control during high volatility periods.

\textbf{Finding}: The hypothesis is \textbf{partially supported}. While LLMs showed better risk-adjusted returns during normal conditions, the maximum drawdown was actually lower for ML models in both markets. This suggests ML models are more conservative in position sizing, while LLMs take on more risk for higher returns.

\subsection{H3: Cost-Efficiency Trade-off}

\textbf{Hypothesis}: The computational overhead of LLMs is justified by improved returns.

\textbf{Finding}: The hypothesis is \textbf{context-dependent}. For small-scale backtesting, the API costs (\$1-5 per backtest) are negligible compared to potential returns. However, for high-frequency applications requiring millisecond latency, ML models remain the only viable option.

\section{Practical Implications}

\subsection{When to Use LLMs}

LLM-based approaches are most suitable for:
\begin{itemize}
    \item Low-frequency trading (daily/weekly signals)
    \item Markets with significant news impact (e.g., emerging markets)
    \item Applications requiring explainability
\end{itemize}

\subsection{When to Use ML}

Traditional ML approaches remain preferable for:
\begin{itemize}
    \item High-frequency trading
    \item Cost-sensitive applications
    \item Markets with limited textual data
\end{itemize}

\section{Limitations and Threats to Validity}

\subsection{Internal Validity}

\begin{itemize}
    \item \textbf{Look-ahead Bias}: While we carefully designed features to avoid future information, subtle biases may remain.
    \item \textbf{Hyperparameter Sensitivity}: Results may vary with different hyperparameter choices.
\end{itemize}

\subsection{External Validity}

\begin{itemize}
    \item \textbf{Market Generalization}: Results are limited to two markets; further validation is needed.
    \item \textbf{Time Period}: The specific time periods studied may not generalize to other market conditions.
\end{itemize}

\subsection{Construct Validity}

\begin{itemize}
    \item \textbf{Evaluation Metrics}: Sharpe ratio and maximum drawdown may not capture all relevant aspects of trading performance.
\end{itemize}

\section{Future Work}

\subsection{Model Improvements}

\begin{itemize}
    \item Fine-tuning LLMs on financial corpora
    \item Hybrid approaches combining ML and LLM strengths
    \item Multi-modal models incorporating charts and images
\end{itemize}

\subsection{Market Expansion}

\begin{itemize}
    \item Validation on additional markets (crypto, commodities)
    \item Cross-market portfolio optimization
\end{itemize}

\subsection{Real-time Implementation}

\begin{itemize}
    \item Latency optimization techniques
    \item Live trading deployment
\end{itemize}